\documentclass[11pt,a4paper]{article}
\usepackage[T1]{fontenc}\usepackage[margin=25mm]{geometry}
\usepackage{amsmath,amssymb,amsthm,mathtools,lmodern,microtype}
\usepackage[hidelinks]{hyperref}\usepackage{xurl}
\hypersetup{pdftitle={A Linear Positivity Region for Thue Morse Pressure},pdfauthor={Research note prepared for Vladimir Reshetnikov with OpenAI}}
\setlength{\parskip}{3pt}
\newtheorem{theorem}{Theorem}[section]\newtheorem{lemma}[theorem]{Lemma}\newtheorem{proposition}[theorem]{Proposition}
\newcommand{\ev}{\operatorname{ev}}\newcommand{\Z}{\mathbb Z}
\title{A Linear Positivity Region for Thue Morse Pressure}
\author{Research note prepared for Vladimir Reshetnikov with OpenAI}\date{1 October 2026}
\begin{document}\maketitle
\begin{abstract}
We prove a positive region of linear width beyond the first nonlinear feedback boundary in integer Thue--Morse pressure. The response and pressure coefficients at offset $s$ are positive whenever $m\ge2048$ and $128s+64\le m$. The proof improves both parts of the preceding estimate: complete multipartite matching counts strengthen the marked-pair denominator, a telescoping trigonometric inequality gives exponential bounds for all auxiliary moments, and a weighted resolvent estimate controls the restored insertion and feedback on a fixed complex disk. The full proposed positive range below degree $6m$ remains open.
\end{abstract}
\paragraph{Status.} This is unrefereed ordinary mathematics with exact arithmetic checks, not a Lean formalization. It is a separate strengthening of the preceding report; earlier artifacts are unchanged. No priority or optimal-constant claim is made.

\section{The linear region}
For $m\ge2$ put $d=2m$ and consider
\[
(V_af)(x)=\sum_{2y=x\bmod1}(\cos\pi y+a\sin\pi y)^d f(y).
\]
Let $h_a(0)=1$ normalize the local analytic eigenfunction, with eigenvalue $\lambda_m(a)$. Define
\[
H_{m,r}=[a^{2r}]h_a(1/2),\qquad E_{m,r}=[t^{2m+2r}]p_m(t/\pi).
\]
Write $R_j=\sum_{w\in\Z+1/2}(\pi w)^{-2j}$, so $H_{m,0}=R_m$ and $0<R_j\le R_1=1$. The exact feedback identities are
\begin{equation}\label{eq:feedback}
\lambda_m(a)=1+a^dh_a(1/2),\qquad p_m(t/\pi)=d\log\cos t+\log\lambda_m(\tan t).
\end{equation}
The preceding triangle and boundary results give $H_{m,r},E_{m,r}>0$ for $1\le r\le m$ \cite{triangle,boundary}. A companion report establishes fixed-offset asymptotics and a logarithmic-width sufficient condition \cite{offset}. Here the offset may be a fixed positive fraction of $m$.
\begin{theorem}\label{thm:main}
For integers $s\ge0$ and $m\ge2048$,
\begin{equation}\label{eq:condition}
128s+64\le m\quad\Longrightarrow\quad H_{m,m+s}>0,\qquad E_{m,m+s}>0.
\end{equation}
Equivalently, $d\ge4096$ and $128(2s+1)\le d$ suffice.
\end{theorem}
Together with the earlier triangle, this proves positivity for all integer indices
\[
1\le r\le m+\left\lfloor\frac{m-64}{128}\right\rfloor,\qquad m\ge2048.
\]
The pressure degree at the upper end is $4m+2\lfloor(m-64)/128\rfloor$. This remains far below the conjectural full range $2m<\deg<6m$. Small-$m$ response positivity cannot extend without qualification: $H_{2,3}=-488/27$ \cite{boundary}.

\section{Stronger marked pairs and tangent moments}
For a summable nonnegative sequence $x$, write
\[
B_x(a)=\prod_i(1+ax_i),\qquad
\mathcal C_{d,q}(x)=[a^q]\prod_i\bigl((1+ax_i)^d-(ax_i)^d\bigr).
\]
The latter polynomial keeps only monomials whose exponent at each coordinate is at most $d-1$.

\begin{lemma}\label{lem:matching}
Let $q=d+2s$, $k=2s+1$, and $d\ge2s+2$. Every color profile of $q$ labeled units with multiplicities at most $d-1$ has at least
\[
(d-1)_k=(d-1)(d-2)\cdots(d-k)
\]
matchings of $k$ disjoint pairs using different colors. Consequently, for $a>0$,
\begin{equation}\label{eq:marked}
\mathcal C_{d,q}(x)\le
\frac{d^{2k}}{k!(d-1)_k}a^{2k-q}e_2(x)^kB_x(a)^d.
\end{equation}
\end{lemma}
\begin{proof}
Let $M_k(\alpha)$ count the matchings for a color vector $\alpha$. We prove by induction on $k$ that moving one unit from a smaller class to a larger class cannot increase $M_k$. The case $k=0$ is constant. For a move $(a,b)\mapsto(a+1,b-1)$ with $a\ge b\ge1$, remove the moved unit. Matchings that leave it unmatched, or pair it with another color, cancel in the difference. Thus
\begin{align*}
&M_k(a+1,b-1,\gamma)-M_k(a,b,\gamma)\\
&\quad=(b-1)M_{k-1}(a,b-2,\gamma)-aM_{k-1}(a-1,b-1,\gamma).
\end{align*}
If $b=1$ the first term is absent. Otherwise the induction hypothesis makes its $M$ value at most the second one. The difference is nonpositive.

Fill the largest class to $d-1$, then merge the remaining units into one class. The remaining total is $k\le d-1$, so all moves respect the cap. The minimizing profile is therefore $(d-1,k)$, whose $k$-matchings number exactly $(d-1)_k$.

For finitely many coordinates introduce independent variables $y_i$ and
\[
\mathcal D_x=\sum_{i<j}x_ix_j\partial_{y_i}\partial_{y_j}.
\]
In $a^{2k}\mathcal D_x^k\prod_i(1+y_i)^d/k!$, evaluated at $y_i=ax_i$, the coefficient of each original monomial is multiplied by its number of unordered labeled-unit matchings. The derivatives choose distinct units, and their $k!$ orders cancel the denominator. At positive $a$, each derivative costs at most $d x_i/(1+ax_i)$. Drop these denominators, sum the pairs, and divide by $a^q(d-1)_k$ to obtain~\eqref{eq:marked}. Monotone convergence gives the infinite-coordinate version.
\end{proof}

Set $\theta_j=\pi/2^{j+1}$ and, for positive odd $n$, $x_j(n)=|\tan(n\theta_j)|/n$. Define
\[
S_s=\sum_{\substack{n\ge3\\n\text{ odd}}}n^{2s}e_2(x(n))^{2s+1}.
\]
\begin{lemma}\label{lem:moment}
For all $s\ge0$,
\begin{equation}\label{eq:moment}
S_s\le(4\pi)^{2s}S_0<11(4\pi)^{2s}.
\end{equation}
\end{lemma}
\begin{proof}
For $0\le a\le1/2$ and real $\theta$,
\begin{equation}\label{eq:telescoping}
(|\cos\theta|+a|\sin\theta|)(1+|\sin\theta|)\le1+|\sin2\theta|.
\end{equation}
Put $u=|\sin\theta|$, $c=|\cos\theta|$. If $u>0$, the inequality is equivalent to
\[
a(1+u)\le c+\frac{u}{1+c}.
\]
The left side is at most one, and the right side is at least $c+u^2/(1+c)=1$. The case $u=0$ is equality.

Multiply~\eqref{eq:telescoping} along $n\theta_j$, $j=1,\ldots,N$. Since $|\sin(n\pi/2)|=1$, the product is at most two. Its infinite limit exists locally uniformly in $a$, with nonnegative coefficients. At $a=1/2$ its degree-two coefficient is at most eight. The dyadic sine-product identity gives
\[
\prod_{j\ge1}|\cos(n\theta_j)|=\frac2{\pi n},
\]
so that coefficient is $2e_2(|\tan(n\theta_j)|)/(\pi n)$. Therefore $n e_2(x(n))\le4\pi$, uniformly in odd $n$. Multiplying this inequality to power $2s$ by $e_2(x(n))$ and summing proves the first part of~\eqref{eq:moment}. The earlier two-insertion estimate gives $S_0<11$ \cite{boundary}.
\end{proof}
The same argument also yields $e_j(|\tan(n\theta_ell)|:\ell\ge1)/n\le\pi2^j$ for every fixed $j$.

\section{A fixed disk for the capped resolvent}
Let $A=V_0$, $h_0=h_{a=0}$, $Pf=f(0)h_0$, $Q=I-P$, and $\ev f=f(1/2)$. Work in the invariant Fourier space with indices $1-m,\ldots,m-1$. Its atomic eigenvalues at zero are $2^{-j}$, $0\le j\le2m-2$, so the eigenvalue one is simple \cite{integer,triangle}.

Write
\[
V_a=\sum_{j=0}^d a^jV_j,\qquad V_{<}(a)=\sum_{j=0}^{d-1}a^jV_j.
\]
Evaluation at zero is preserved by $V_{<}$. Let $f_{<}(a)$ be its normalized fixed vector, and put
\[
T_{<}(a)=(I-QV_{<}(a)Q)^{-1}\ \text{on }Q,\qquad Z(a)=\ev f_{<}(a).
\]
For $x\in[0,1]$ define, periodically,
\[
g(x)=\sin\pi x,\quad H(x)=1+\tfrac18\sin\pi x,\quad K_d(x)=H(x)^dg(x).
\]
Use $\|f\|_K=\sup|f|/K_d$ on vectors vanishing at zero.

\begin{lemma}\label{lem:disk}
For $d\ge128$ and complex $|a|\le1/16$, the restriction of $V_{<}(a)$ to $Q$ has $K$-norm at most $3/4$. Hence $\|T_{<}(a)\|_K\le4$ throughout this disk.
\end{lemma}
\begin{proof}
Put $a_0=1/16$, $\beta=1/8$, $u=\sin(\pi x/2)$, $c=\cos(\pi x/2)$, and $N=1+2\beta uc$. Define
\[
p_1=\frac{(c+a_0u)(1+\beta u)}N,\quad r_1=\frac{a_0u(1+\beta u)}N,
\]
\[
p_2=\frac{(u+a_0c)(1+\beta c)}N,\quad r_2=\frac{a_0c(1+\beta c)}N.
\]
The absolute capped operator applied to $K_d$, divided by $K_d(x)$, is
\begin{equation}\label{eq:branches}
\frac{p_1^d-r_1^d}{2c}+\frac{p_2^d-r_2^d}{2u},
\end{equation}
with endpoint limits. This majorizes complex $a$ because the sum of absolute coefficients at radius $a_0$ is $(|\cos\pi y|+a_0\sin\pi y)^d-(a_0\sin\pi y)^d$.

Both $p_i\le1$. Indeed,
\[
N-(c+a_0u)(1+2a_0u)=\frac{u^2}{1+c}+a_0u(2c-1)-2a_0^2u^2.
\]
If $c\ge1/2$, this is at least $u^2(1/2-2a_0^2)$. If $c\le1/2$, it is at least $1/2-a_0-2a_0^2=55/128$.

When $u\le1/4$, the first term of~\eqref{eq:branches} is at most $2/\sqrt{15}<8/15$. The second is at most
\[
\frac{9d}{16}(45/128)^{d-1},
\]
since $p_2\le45/128$ and $p_2-r_2=u(1+\beta c)/N$. For $d\ge4$ this decreases with $d$, and at $d=4$ its sum with $8/15$ is below $3/4$. The case $c\le1/4$ is symmetric.

In the remaining interior, each $p_i\le39/40$. If $c\ge1/2$, its normalized numerator gap is at least $(63/2048)/(9/8)=7/256>1/40$; if $c\le1/2$ it is even larger. Thus~\eqref{eq:branches} is at most $4(39/40)^d\le4(39/40)^{128}<3/4$. The Neumann series on $Q$ completes the proof.
\end{proof}

Put
\[
L=\frac{153}{128},\qquad M=\frac{1377}{1024},\qquad J=L^2=\frac{23409}{16384}.
\]
The elementary insertion bounds, with $1\le j<d$, are
\[
Ag\le g/2,\quad |V_jf|\le\tfrac12\binom dj\|f\|_\infty g,\quad\|h_0\|_\infty\le1.
\]
They follow from the two branches and weighted AM--GM as in \cite{boundary,offset}. Since $K_d\ge g$, Lemma~\ref{lem:disk} gives
\[
\|f_{<}-h_0\|_K\le2(17/16)^d,\qquad\|f_{<}\|_\infty\le3L^d.
\]
The initial $h_0$ is used in sup norm; it does not vanish at zero.

For a trigonometric polynomial $v$ of degree at most $m-1$, Fourier coefficient bounds give $\|v'\|_\infty\le2\pi m(m-1)\|v\|_\infty$. Since $QV_df$ vanishes at zero and $\|QV_df\|_\infty\le2\|f\|_\infty$,
\[
\|QV_df\|_K\le\|QV_df\|_g<2d^2\|f\|_\infty.
\]
Therefore, throughout $|a|\le1/16$,
\begin{equation}\label{eq:analyticbounds}
|\ev T_{<}QV_df_{<}|\le24d^2M^d,\qquad
|Z\,\ev T_{<}(f_{<}-h_0)|\le24J^d.
\end{equation}
For the second bound, $\|T_{<}(f_{<}-h_0)\|_\infty\le8L^d$.

\section{The exact response comparison}
The finite-iterate periodization yields the absolutely convergent capped orbit formula
\begin{equation}\label{eq:orbit}
[a^q]Z(a)=\sum_{w\in\Z+1/2}(\pi w)^{-d}
[a^q]\prod_{j\ge1}\bigl((1+a\tan(\pi w/2^j))^d-(a\tan(\pi w/2^j))^d\bigr).
\end{equation}
At every fixed coefficient all insertion degrees are below $d$. Their absolute values map bounded functions into multiples of $g$, and all zero-degree gaps contract $g$ by $1/2$. Summing over the finitely many degree compositions proves absolute convergence and the passage from finite iterates \cite{offset}.

Let $n=2s<d$ and $q=d+n$. Applying $Q$ to the true eigenvalue equation and subtracting the capped fixed-vector equation gives
\[
h_a-f_{<}=a^dT_{<}\bigl(QV_dh_a-h_a(1/2)(h_a-h_0)\bigr).
\]
At total order $d+n<2d$, replace $h_a$ in the braces by $f_{<}$. Thus
\begin{equation}\label{eq:restore}
H_{m,m+s}=[a^q]Z+[a^n]\bigl(\ev T_{<}QV_df_{<}-Z\,\ev T_{<}(f_{<}-h_0)\bigr).
\end{equation}
Cauchy's inequality and $M<J$ bound the second term in absolute value by
\begin{equation}\label{eq:restorebound}
24(d^2+1)J^d16^n.
\end{equation}
This is the key improvement over low-order polynomial-in-$d$ jet bounds. No sign of the frozen response is assumed.

\section{Three exponentially small errors}
Let $b_j=\tan(\pi/2^{j+1})$ and $B(a)=\prod_j(1+ab_j)$. The strict tangent comparison from \cite{boundary} is
\[
B_{x(n)}(a)\le\eta(a)B(a),\qquad\eta(a)=\frac{1+\theta a}{1+a},\qquad\theta=\frac{1+\sqrt2}{3}<\frac{81}{100}.
\]
It follows from weak log-majorization by $(\theta,b_2,b_3,\ldots)$; the earlier report supplies the dyadic single-crossing proof.

Assume $d\ge4096$ and $k=2s+1\le d/128$. Choose $a_q$ with $a_qB'(a_q)/B(a_q)=q/d$. The exact half-angle certificate gives $a_q>4/5$ \cite{boundary}. For the upper bound, $q/d<33/32$ and
\[
B'(1)/B(1)>\tfrac12+\tfrac3{11}+\tfrac3{19}+\tfrac3{35}+\tfrac3{67}>33/32,
\]
using $\tan x>x$ and $\pi>3$. Thus $a_q<1$. Consequently
\[
\eta(a_q)<206/225,\qquad B(a_q)/a_q^{q/d}>3.
\]
Indeed, using $b_1=1,b_2>2/5,b_3>1/6$,
\[
\frac{B(a_q)}{a_q^{q/d}}\ge\frac{B(a_q)}{a_q}
>\frac1{a_q}+\frac{47}{30}+\frac{19a_q}{30}+\frac{a_q^2}{15}>3.
\] The $a_q$-tilted countable Bernoulli sum has integer mean $q$ and variance at most $q$. The mean is a mode by \cite[Theorem 2]{gnedin}; Chebyshev's inequality then gives
\begin{equation}\label{eq:mainlower}
E_0:=2(2/\pi)^d[u^q]B(u)^d>\frac{(21/11)^d}{4\sqrt q}.
\end{equation}
Here $E_0$ is the unmodified nearest-orbit contribution. We used $[u^q]B(u)^d\ge B(a_q)^da_q^{-q}/(8\sqrt q)$ and $\pi<22/7$.

\paragraph{Remote orbits.}
Lemmas~\ref{lem:matching} and~\ref{lem:moment} bound their absolute sum divided by $E_0$ by
\[
8\sqrt q\,\frac{d^{2k}S_s}{k!(d-1)_k}\eta(a_q)^d.
\]
Use $S_s<11\cdot13^k$, $(d-1)_k\ge(127d/128)^k$, and $k!\ge(k/3)^k$. For $v=k/d\le1/128$, the increasing function
\[
v\log\frac{39\cdot128}{127v}
\]
is at most $\log(39\cdot128^2/127)/128<9/128$. Hence the remote ratio is at most
\[
132\sqrt d\,\exp\bigl((9/128-19/225)d\bigr).
\]
Its logarithm divided by $d$ is bounded by
\[
\frac{19}{8192}+\frac9{128}-\frac{19}{225}<-\frac1{100}.
\]
Here and below we use $d\ge4096$, $\log4096<9$, and the decreasing function $(c+b\log d)/d$.

\paragraph{Deleted nearest monomials.}
Since $q<2d$, deleting full powers from the nearest product removes exactly
\[
\sum_jb_j^d[u^n]\prod_{i\ne j}(1+ub_i)^d\le\frac{2(2d)^n}{n!}.
\]
The relative cost is at most $16\sqrt q(2d)^n/(n!3^d)$. For $n>0$, use $n!\ge(n/3)^n$ and $n/d\le1/128$; for $n=0$ the numerator factor is one. The normalized logarithm is at most
\[
\frac{17}{8192}+\frac8{128}-1<-\frac12.
\]

\paragraph{Restoration and feedback.}
Divide~\eqref{eq:restorebound} by~\eqref{eq:mainlower}. The rational inequality
\[
J\frac{11}{21}=\frac{257499}{344064}<\frac34
\]
gives the relative bound $288d^{5/2}(3/4)^d16^n$. Its normalized logarithm is at most
\[
\frac{57}{8192}-\frac{29}{128}<-\frac15,
\]
using $\log16<3$ and $\log(4/3)>1/4$.

Each of these three errors is less than $E_0/8$. Thus
\begin{equation}\label{eq:positiveH}
H_{m,m+s}>\frac{E_0}{2}>\frac{(21/11)^d}{8\sqrt q}>0.
\end{equation}
The factorial bound follows from $\log(k!)\ge\int_1^k\log x\,dx\ge k\log k-k$ and $e<3$. All displayed rational margins and elementary exponential bounds are checked exactly in the source bundle.

\section{Pressure transfer and scope}
Every smaller offset satisfies the same linear condition. The pre-feedback triangle and~\eqref{eq:positiveH} make all response terms in
\begin{align*}
E_{m,m+s}={}&\sum_{j=0}^{m+s}H_{m,j}[t^{2(m+s-j)}](\tan t/t)^{d+2j}\\
&-\tfrac12[t^n](\tan t/t)^{2d}h_{\tan t}(1/2)^2
-\frac{m}{2m+s}R_{2m+s}
\end{align*}
positive, apart from the two explicit subtractions. This formula follows from~\eqref{eq:feedback}; the cubic logarithm begins at degree $3d>2d+n$.

Take $|t|=\rho=1/32$. Positive tangent coefficients give
\[
|\tan t|\le\tan\rho\le\rho/(1-\rho^2)<1/16.
\]
The coefficients through degree $n<d$ of the true response equal those of $f_{<}$. Cauchy's inequality and~\eqref{eq:analyticbounds} bound the quadratic subtraction by
\[
\tfrac92\bigl[J(1024/1023)^2\bigr]^d32^n.
\]
The cosine subtraction is at most one, so their sum is bounded by the same expression with $9/2$ replaced by $11/2$. The exact inequality
\[
J(1024/1023)^2\frac{11}{21}<\frac34
\]
then bounds its ratio to~\eqref{eq:positiveH} by $66\sqrt d(3/4)^d32^n$. Its normalized logarithm is at most
\[
\frac{19}{8192}-\frac{113}{512}<0,
\]
using $\log32<15/4$. The single response term $H_{m,m+s}$ therefore exceeds both subtractions, completing Theorem~\ref{thm:main}.

The bundle includes exact checks of the strengthened matching multiplicity on $2,632$ finite profiles, all branch-contraction and exponential-margin constants, and the earlier response identity checks. These supplement the universal analytic proof; no negative finite search is used as a proof.

The remaining pressure interval up to degree $6m$, the optimal linear-width constant, and the first later negative pressure degree remain open. The fixed-saddle asymptotic of \cite{offset} is not asserted uniformly in this linear region; a moving saddle would be required for such a statement.

\begin{thebibliography}{9}\small\raggedright
\bibitem{integer} V. Reshetnikov, \emph{Integer Thue--Morse Pressure}, ProveIt incoming source, revision \texttt{6a98f89bac85cc6b4ba5d5c3b63996037d9824e1}, blob \texttt{1973fd17245864bb8a83e1b1d90416ac032ba7c1}.
\bibitem{triangle} \emph{The Full Pre Feedback Positivity Triangle for Thue Morse Pressure}, companion research report prepared for V. Reshetnikov with OpenAI, 1 October 2026.
\bibitem{boundary} \emph{Positivity at the First Feedback Boundary in Thue Morse Pressure}, companion research report prepared for V. Reshetnikov with OpenAI, 1 October 2026.
\bibitem{offset} \emph{Positivity Beyond the First Thue Morse Feedback Boundary}, companion research report prepared for V. Reshetnikov with OpenAI, 1 October 2026.
\bibitem{gnedin} A. Gnedin, \emph{Cross Modality of the Extended Binomial Sums}, arXiv:2408.06477v1 (2024), Theorem 2. \url{https://arxiv.org/html/2408.06477v1}.
\end{thebibliography}
\end{document}
